% Fragmento público generado desde propietarios íntegros.
% Residencia: 52.8.1.
% Fuente: ../incoming/radion_monografia_20260827/manuscrito/sections/09_sector_90120.tex:176-244
\subsubsection{Jerarquía armónica global y semigrupo \(\langle90,120\rangle\)}

Para evitar una colisión de signos, fijamos
\[
R_{\mathrm{tow}}:=-R_{\mathrm{ch}}.
\]
Entonces
\begin{equation}
T=e^{-xI+yR_{\mathrm{tow}}}
=q_-P_+^{\mathrm{tow}}+q_+P_-^{\mathrm{tow}}.
\label{eq:tower-operator}
\end{equation}
Las torres pares e impares son
\begin{align}
c_n&=\operatorname{Tr}(T^n)
=q_-^n+q_+^n
=2e^{-nx}\cosh(ny),
\label{eq:cn}\\
s_n&=\operatorname{Tr}(R_{\mathrm{tow}}T^n)
=q_-^n-q_+^n
=2e^{-nx}\sinh(ny).
\label{eq:sn}
\end{align}
Bajo \(C^*\mapsto-C^*\), \(c_n\) es par y \(s_n\) es impar. Se cumple
\begin{equation}
c_n^2-s_n^2=4e^{-2nx},
\label{eq:tower-invariant}
\end{equation}
y ambas sucesiones satisfacen la recurrencia de orden dos determinada por los
autovalores \(q_\pm\).

El semigrupo
\begin{equation}
\langle90,120\rangle=30\langle3,4\rangle
\label{eq:semigroup}
\end{equation}
es la jerarquía armónica de un único operador bicapa. No es una lista de ocho
correcciones laterales. Su completitud significa que toda razón positiva
admite una expansión convergente; no significa que una expansión obtenida a
partir del blanco sea un selector físico prospectivo.

\subsubsection{Clasificación tipada y separación de las realizaciones \(90/120\)}

\begin{longtable}{>{\bfseries}p{.19\textwidth} p{.24\textwidth} p{.45\textwidth}}
\toprule
Objeto & Operación & Qué conserva / por qué no es otro objeto\\
\midrule
\endfirsthead
\toprule
Objeto & Operación & Qué conserva / por qué no es otro objeto\\
\midrule
\endhead
Incidencia \(N_6,N_8\) & conteo finito sobre fases activas & cardinalidad y orientación; produce \(90,120\), no una serie;\\
\(-1/12\) & diferencia normalizada \eqref{eq:minus-one-twelve} & escalar de incidencia; no peso Barbero;\\
\(E_{\mathrm{dod}}\) & extractor armónico de la rueda & fase dodecafásica; no acarreo real;\\
\(h=C^*/6\) & reconstrucción de seis parejas & apertura por pareja; no normalización metrológica;\\
\(F_{\mathrm{spec}}\) & cola logarítmica \(j\ge2\) & publicación espectral completa; difiere de \(\epsR^\circ\);\\
\(\epsone\) & resta APP entre \(S_T\) y \(S_E\) & empalme real de dos lifts; no serie de Lambert;\\
\(K_{6\to8}\) & \(12\Delta S_{90}-\Delta S_{120}\) & salto marcado y evaluación conjugada; coeficiente de polo \(1/24\) comprobado; salida Barbero;\\
\(12:-1\) & doce sectores y una costura orientada & peso del núcleo; certificado posterior \eqref{eq:pole-condition}; distinto de \(-1/12\);\\
\(\langle90,120\rangle\) & semigrupo de potencias de \(T\) & jerarquía armónica global; no residuo por especie.\\
\bottomrule
\end{longtable}

\begin{thesisbox}
Todos estos objetos descienden del mismo estado APP--TRIT--TPK. Esa genealogía
común explica por qué comparten números y simetrías. Su tipado explica por qué
no pueden reemplazarse mutuamente.
\end{thesisbox}
